
\documentclass{article}
\usepackage{colortbl}
\usepackage{makecell}
\usepackage{multirow}
\usepackage{supertabular}

\begin{document}

\newcounter{utterance}

\centering \large Interaction Transcript for game `matchit\_ascii', experiment `similar\_grid\_1', episode 3 with deepseek{-}v4{-}pro.
\vspace{24pt}

{ \footnotesize  \setcounter{utterance}{1}
\setlength{\tabcolsep}{0pt}
\begin{supertabular}{c@{$\;$}|p{.15\linewidth}@{}p{.15\linewidth}p{.15\linewidth}p{.15\linewidth}p{.15\linewidth}p{.15\linewidth}}
    \# & $\;$A & \multicolumn{4}{c}{Game Master} & $\;\:$B\\
    \hline

    \theutterance \stepcounter{utterance}  
    & & \multicolumn{4}{p{0.6\linewidth}}{
        \cellcolor[rgb]{0.9,0.9,0.9}{
            \makecell[{{p{\linewidth}}}]{
                \texttt{\tiny{[P1$\langle$GM]}}
                \texttt{Συμμετέχετε σε ένα συνεργατικό παιχνίδι μαντεψιάς.} \\
\texttt{Ο στόχος είναι να βρείτε αν αυτό το πλέγμα και μια άλλη εικόνα που μόνο εγώ μπορώ να δω, είναι τα ίδια.} \\
\texttt{Παρακαλώ περιγράψτε πρώτα το πλέγμα σας. Στη συνέχεια, θα δώσω τη δική μου περιγραφή και μπορούμε να κάνουμε ο ένας στον άλλο ερωτήσεις σχετικά με τα πλέγματα για να καταλάβουμε αν είναι τα ίδια. Ξεκινήστε τώρα τη σύντομη περιγραφή του πλέγματός σας με "ΠΕΡΙΓΡΑΦΗ:" ακολουθούμενη από την περιγραφή. Μην προσθέσετε τίποτα άλλο. Αυτό είναι το πλέγμα σας:} \\
\texttt{X X X X X} \\
\texttt{X ▢ ▢ ▢ ▢} \\
\texttt{X X X ▢ ▢} \\
\texttt{X ▢ ▢ ▢ ▢} \\
\texttt{X ▢ ▢ ▢ ▢} \\
            }
        }
    }
    & \\ \\

    \theutterance \stepcounter{utterance}  
    & \multicolumn{4}{p{0.6\linewidth}}{
        \cellcolor[rgb]{0.9,0.9,0.9}{
            \makecell[{{p{\linewidth}}}]{
                \texttt{\tiny{[P1$\rangle$GM]}}
                \texttt{ΠΕΡΙΓΡΑΦΗ: Ένα πλέγμα 5x5. Η πρώτη στήλη είναι γεμάτη με X. Η πρώτη γραμμή είναι γεμάτη με X. Η τρίτη γραμμή έχει X στις τρεις πρώτες στήλες και τα υπόλοιπα δύο είναι κενά. Όλα τα υπόλοιπα κελιά είναι κενά.} \\
            }
        }
    }
    & & \\ \\

    \theutterance \stepcounter{utterance}  
    & & & \multicolumn{2}{p{0.3\linewidth}}{
        \cellcolor[rgb]{0.9,0.9,0.9}{
            \makecell[{{p{\linewidth}}}]{
                \texttt{\tiny{[GM$|$GM]}}
                \texttt{continue} \\
            }
        }
    }
    & & \\ \\

    \theutterance \stepcounter{utterance}  
    & & \multicolumn{4}{p{0.6\linewidth}}{
        \cellcolor[rgb]{0.9,0.9,0.9}{
            \makecell[{{p{\linewidth}}}]{
                \texttt{\tiny{[GM$\rangle$P2]}}
                \texttt{Συμμετέχετε σε ένα συνεργατικό παιχνίδι μαντεψιάς.} \\
\texttt{Ο στόχος είναι να βρείτε αν αυτό το πλέγμα και μια άλλη εικόνα που μόνο εγώ μπορώ να δω, είναι τα ίδια.} \\
\texttt{Παρακαλώ περιγράψτε πρώτα το πλέγμα σας. Στη συνέχεια, θα δώσω τη δική μου περιγραφή και μπορούμε να κάνουμε ο ένας στον άλλο ερωτήσεις σχετικά με τα πλέγματα για να καταλάβουμε αν είναι τα ίδια. Ξεκινήστε τώρα τη σύντομη περιγραφή του πλέγματός σας με "ΠΕΡΙΓΡΑΦΗ:" ακολουθούμενη από την περιγραφή. Μην προσθέσετε τίποτα άλλο. Αυτό είναι το πλέγμα σας:} \\
\texttt{X ▢ ▢ ▢ ▢} \\
\texttt{X ▢ ▢ ▢ ▢} \\
\texttt{X X X ▢ ▢} \\
\texttt{X ▢ ▢ ▢ ▢} \\
\texttt{X X X X X} \\
            }
        }
    }
    & \\ \\

    \theutterance \stepcounter{utterance}  
    & & & \multicolumn{4}{p{0.6\linewidth}}{
        \cellcolor[rgb]{0.9,0.9,0.9}{
            \makecell[{{p{\linewidth}}}]{
                \texttt{\tiny{[GM$\langle$P2]}}
                \texttt{ΠΕΡΙΓΡΑΦΗ: Πλέγμα 5x5. Πρώτη στήλη γεμάτη X. Τρίτη γραμμή έχει τρία X στην αρχή. Τελευταία γραμμή γεμάτη X. Οι υπόλοιπες θέσεις είναι κενές (▢).} \\
            }
        }
    }
     \\ \\

    \theutterance \stepcounter{utterance}  
    & & & \multicolumn{2}{p{0.3\linewidth}}{
        \cellcolor[rgb]{0.9,0.9,0.9}{
            \makecell[{{p{\linewidth}}}]{
                \texttt{\tiny{[GM$|$GM]}}
                \texttt{continue} \\
            }
        }
    }
    & & \\ \\

    \theutterance \stepcounter{utterance}  
    & & \multicolumn{4}{p{0.6\linewidth}}{
        \cellcolor[rgb]{0.9,0.9,0.9}{
            \makecell[{{p{\linewidth}}}]{
                \texttt{\tiny{[GM$\rangle$P2]}}
                \texttt{Αυτή είναι η ΠΕΡΙΓΡΑΦΗ μουΠΕΡΙΓΡΑΦΗ: Ένα πλέγμα 5x5. Η πρώτη στήλη είναι γεμάτη με X. Η πρώτη γραμμή είναι γεμάτη με X. Η τρίτη γραμμή έχει X στις τρεις πρώτες στήλες και τα υπόλοιπα δύο είναι κενά. Όλα τα υπόλοιπα κελιά είναι κενά.} \\
\texttt{Τώρα κάντε μια ερώτηση για να μάθετε νέες πτυχές του πλέγματός μου που μπορεί να διαφέρουν από το δικό σας πλέγμα. Ξεκινήστε με "ΕΡΩΤΗΣΗ:" και μην προσθέσετε τίποτα άλλο.} \\
            }
        }
    }
    & \\ \\

    \theutterance \stepcounter{utterance}  
    & & & \multicolumn{4}{p{0.6\linewidth}}{
        \cellcolor[rgb]{0.9,0.9,0.9}{
            \makecell[{{p{\linewidth}}}]{
                \texttt{\tiny{[GM$\langle$P2]}}
                \texttt{ΕΡΩΤΗΣΗ: Στο πλέγμα σας, η τελευταία (5η) γραμμή έχει όλα τα κελιά γεμάτα με X;} \\
            }
        }
    }
     \\ \\

    \theutterance \stepcounter{utterance}  
    & & & \multicolumn{2}{p{0.3\linewidth}}{
        \cellcolor[rgb]{0.9,0.9,0.9}{
            \makecell[{{p{\linewidth}}}]{
                \texttt{\tiny{[GM$|$GM]}}
                \texttt{continue} \\
            }
        }
    }
    & & \\ \\

    \theutterance \stepcounter{utterance}  
    & & \multicolumn{4}{p{0.6\linewidth}}{
        \cellcolor[rgb]{0.9,0.9,0.9}{
            \makecell[{{p{\linewidth}}}]{
                \texttt{\tiny{[P1$\langle$GM]}}
                \texttt{Αυτή είναι η ΠΕΡΙΓΡΑΦΗ μουΠΕΡΙΓΡΑΦΗ: Πλέγμα 5x5. Πρώτη στήλη γεμάτη X. Τρίτη γραμμή έχει τρία X στην αρχή. Τελευταία γραμμή γεμάτη X. Οι υπόλοιπες θέσεις είναι κενές (▢).} \\
\texttt{ΕΡΩΤΗΣΗ: Στο πλέγμα σας, η τελευταία (5η) γραμμή έχει όλα τα κελιά γεμάτα με X;Ξεκινήστε την απάντησή σας με "ΑΠΑΝΤΗΣΗ:" και μην προσθέσετε τίποτα άλλο.} \\
            }
        }
    }
    & \\ \\

    \theutterance \stepcounter{utterance}  
    & \multicolumn{4}{p{0.6\linewidth}}{
        \cellcolor[rgb]{0.9,0.9,0.9}{
            \makecell[{{p{\linewidth}}}]{
                \texttt{\tiny{[P1$\rangle$GM]}}
                \texttt{ΑΠΑΝΤΗΣΗ: Όχι, η τελευταία γραμμή έχει X μόνο στην πρώτη στήλη και τα υπόλοιπα είναι κενά.} \\
            }
        }
    }
    & & \\ \\

    \theutterance \stepcounter{utterance}  
    & & & \multicolumn{2}{p{0.3\linewidth}}{
        \cellcolor[rgb]{0.9,0.9,0.9}{
            \makecell[{{p{\linewidth}}}]{
                \texttt{\tiny{[GM$|$GM]}}
                \texttt{continue} \\
            }
        }
    }
    & & \\ \\

    \theutterance \stepcounter{utterance}  
    & & \multicolumn{4}{p{0.6\linewidth}}{
        \cellcolor[rgb]{0.9,0.9,0.9}{
            \makecell[{{p{\linewidth}}}]{
                \texttt{\tiny{[P1$\langle$GM]}}
                \texttt{Τώρα κάντε μια ερώτηση για να μάθετε νέες πτυχές του πλέγματός μου που μπορεί να διαφέρουν από το δικό σας πλέγμα. Ξεκινήστε με "ΕΡΩΤΗΣΗ:" και μην προσθέσετε τίποτα άλλο.} \\
            }
        }
    }
    & \\ \\

    \theutterance \stepcounter{utterance}  
    & \multicolumn{4}{p{0.6\linewidth}}{
        \cellcolor[rgb]{0.9,0.9,0.9}{
            \makecell[{{p{\linewidth}}}]{
                \texttt{\tiny{[P1$\rangle$GM]}}
                \texttt{ΕΡΩΤΗΣΗ: Η δεύτερη γραμμή του πλέγματός σας έχει X μόνο στην πρώτη στήλη και όλα τα υπόλοιπα είναι κενά;} \\
            }
        }
    }
    & & \\ \\

    \theutterance \stepcounter{utterance}  
    & & & \multicolumn{2}{p{0.3\linewidth}}{
        \cellcolor[rgb]{0.9,0.9,0.9}{
            \makecell[{{p{\linewidth}}}]{
                \texttt{\tiny{[GM$|$GM]}}
                \texttt{continue} \\
            }
        }
    }
    & & \\ \\

    \theutterance \stepcounter{utterance}  
    & & \multicolumn{4}{p{0.6\linewidth}}{
        \cellcolor[rgb]{0.9,0.9,0.9}{
            \makecell[{{p{\linewidth}}}]{
                \texttt{\tiny{[GM$\rangle$P2]}}
                \texttt{ΑΠΑΝΤΗΣΗ: Όχι, η τελευταία γραμμή έχει X μόνο στην πρώτη στήλη και τα υπόλοιπα είναι κενά.} \\
\texttt{ΕΡΩΤΗΣΗ: Η δεύτερη γραμμή του πλέγματός σας έχει X μόνο στην πρώτη στήλη και όλα τα υπόλοιπα είναι κενά;Ξεκινήστε την απάντησή σας με "ΑΠΑΝΤΗΣΗ:" και μην προσθέσετε τίποτα άλλο.} \\
            }
        }
    }
    & \\ \\

    \theutterance \stepcounter{utterance}  
    & & & \multicolumn{4}{p{0.6\linewidth}}{
        \cellcolor[rgb]{0.9,0.9,0.9}{
            \makecell[{{p{\linewidth}}}]{
                \texttt{\tiny{[GM$\langle$P2]}}
                \texttt{ΑΠΑΝΤΗΣΗ: Ναι, η δεύτερη σειρά έχει X μόνο στην πρώτη στήλη και τα υπόλοιπα είναι κενά.} \\
            }
        }
    }
     \\ \\

    \theutterance \stepcounter{utterance}  
    & & & \multicolumn{2}{p{0.3\linewidth}}{
        \cellcolor[rgb]{0.9,0.9,0.9}{
            \makecell[{{p{\linewidth}}}]{
                \texttt{\tiny{[GM$|$GM]}}
                \texttt{continue} \\
            }
        }
    }
    & & \\ \\

    \theutterance \stepcounter{utterance}  
    & & \multicolumn{4}{p{0.6\linewidth}}{
        \cellcolor[rgb]{0.9,0.9,0.9}{
            \makecell[{{p{\linewidth}}}]{
                \texttt{\tiny{[GM$\rangle$P2]}}
                \texttt{Τώρα κάντε μια ερώτηση για να μάθετε νέες πτυχές του πλέγματός μου που μπορεί να διαφέρουν από το δικό σας πλέγμα. Ξεκινήστε με "ΕΡΩΤΗΣΗ:" και μην προσθέσετε τίποτα άλλο.} \\
            }
        }
    }
    & \\ \\

    \theutterance \stepcounter{utterance}  
    & & & \multicolumn{4}{p{0.6\linewidth}}{
        \cellcolor[rgb]{0.9,0.9,0.9}{
            \makecell[{{p{\linewidth}}}]{
                \texttt{\tiny{[GM$\langle$P2]}}
                \texttt{ΕΡΩΤΗΣΗ: Η τέταρτη γραμμή του πλέγματός σας έχει X μόνο στην πρώτη στήλη και τα υπόλοιπα κελιά είναι κενά;} \\
            }
        }
    }
     \\ \\

    \theutterance \stepcounter{utterance}  
    & & & \multicolumn{2}{p{0.3\linewidth}}{
        \cellcolor[rgb]{0.9,0.9,0.9}{
            \makecell[{{p{\linewidth}}}]{
                \texttt{\tiny{[GM$|$GM]}}
                \texttt{continue} \\
            }
        }
    }
    & & \\ \\

    \theutterance \stepcounter{utterance}  
    & & \multicolumn{4}{p{0.6\linewidth}}{
        \cellcolor[rgb]{0.9,0.9,0.9}{
            \makecell[{{p{\linewidth}}}]{
                \texttt{\tiny{[P1$\langle$GM]}}
                \texttt{ΑΠΑΝΤΗΣΗ: Ναι, η δεύτερη σειρά έχει X μόνο στην πρώτη στήλη και τα υπόλοιπα είναι κενά.} \\
\texttt{ΕΡΩΤΗΣΗ: Η τέταρτη γραμμή του πλέγματός σας έχει X μόνο στην πρώτη στήλη και τα υπόλοιπα κελιά είναι κενά;Ξεκινήστε την απάντησή σας με "ΑΠΑΝΤΗΣΗ:" και μην προσθέσετε τίποτα άλλο.} \\
            }
        }
    }
    & \\ \\

    \theutterance \stepcounter{utterance}  
    & \multicolumn{4}{p{0.6\linewidth}}{
        \cellcolor[rgb]{0.9,0.9,0.9}{
            \makecell[{{p{\linewidth}}}]{
                \texttt{\tiny{[P1$\rangle$GM]}}
                \texttt{ΑΠΑΝΤΗΣΗ: Ναι, η τέταρτη γραμμή έχει X μόνο στην πρώτη στήλη και τα υπόλοιπα είναι κενά.} \\
            }
        }
    }
    & & \\ \\

    \theutterance \stepcounter{utterance}  
    & & & \multicolumn{2}{p{0.3\linewidth}}{
        \cellcolor[rgb]{0.9,0.9,0.9}{
            \makecell[{{p{\linewidth}}}]{
                \texttt{\tiny{[GM$|$GM]}}
                \texttt{continue} \\
            }
        }
    }
    & & \\ \\

    \theutterance \stepcounter{utterance}  
    & & \multicolumn{4}{p{0.6\linewidth}}{
        \cellcolor[rgb]{0.9,0.9,0.9}{
            \makecell[{{p{\linewidth}}}]{
                \texttt{\tiny{[P1$\langle$GM]}}
                \texttt{Τώρα κάντε μια ερώτηση για να μάθετε νέες πτυχές του πλέγματός μου που μπορεί να διαφέρουν από το δικό σας πλέγμα. Ξεκινήστε με "ΕΡΩΤΗΣΗ:" και μην προσθέσετε τίποτα άλλο.} \\
            }
        }
    }
    & \\ \\

    \theutterance \stepcounter{utterance}  
    & \multicolumn{4}{p{0.6\linewidth}}{
        \cellcolor[rgb]{0.9,0.9,0.9}{
            \makecell[{{p{\linewidth}}}]{
                \texttt{\tiny{[P1$\rangle$GM]}}
                \texttt{ΕΡΩΤΗΣΗ: Στο πλέγμα σας, η πρώτη γραμμή έχει X μόνο στην πρώτη στήλη ή είναι γεμάτη με X σε όλες τις στήλες;} \\
            }
        }
    }
    & & \\ \\

    \theutterance \stepcounter{utterance}  
    & & & \multicolumn{2}{p{0.3\linewidth}}{
        \cellcolor[rgb]{0.9,0.9,0.9}{
            \makecell[{{p{\linewidth}}}]{
                \texttt{\tiny{[GM$|$GM]}}
                \texttt{continue} \\
            }
        }
    }
    & & \\ \\

    \theutterance \stepcounter{utterance}  
    & & \multicolumn{4}{p{0.6\linewidth}}{
        \cellcolor[rgb]{0.9,0.9,0.9}{
            \makecell[{{p{\linewidth}}}]{
                \texttt{\tiny{[GM$\rangle$P2]}}
                \texttt{ΑΠΑΝΤΗΣΗ: Ναι, η τέταρτη γραμμή έχει X μόνο στην πρώτη στήλη και τα υπόλοιπα είναι κενά.} \\
\texttt{ΕΡΩΤΗΣΗ: Στο πλέγμα σας, η πρώτη γραμμή έχει X μόνο στην πρώτη στήλη ή είναι γεμάτη με X σε όλες τις στήλες;Ξεκινήστε την απάντησή σας με "ΑΠΑΝΤΗΣΗ:" και μην προσθέσετε τίποτα άλλο.} \\
            }
        }
    }
    & \\ \\

    \theutterance \stepcounter{utterance}  
    & & & \multicolumn{4}{p{0.6\linewidth}}{
        \cellcolor[rgb]{0.9,0.9,0.9}{
            \makecell[{{p{\linewidth}}}]{
                \texttt{\tiny{[GM$\langle$P2]}}
                \texttt{ΑΠΑΝΤΗΣΗ: Η πρώτη γραμμή έχει X μόνο στην πρώτη στήλη, τα υπόλοιπα είναι κενά.} \\
            }
        }
    }
     \\ \\

    \theutterance \stepcounter{utterance}  
    & & & \multicolumn{2}{p{0.3\linewidth}}{
        \cellcolor[rgb]{0.9,0.9,0.9}{
            \makecell[{{p{\linewidth}}}]{
                \texttt{\tiny{[GM$|$GM]}}
                \texttt{continue} \\
            }
        }
    }
    & & \\ \\

    \theutterance \stepcounter{utterance}  
    & & \multicolumn{4}{p{0.6\linewidth}}{
        \cellcolor[rgb]{0.9,0.9,0.9}{
            \makecell[{{p{\linewidth}}}]{
                \texttt{\tiny{[GM$\rangle$P2]}}
                \texttt{Τώρα κάντε μια ερώτηση για να μάθετε νέες πτυχές του πλέγματός μου που μπορεί να διαφέρουν από το δικό σας πλέγμα. Ξεκινήστε με "ΕΡΩΤΗΣΗ:" και μην προσθέσετε τίποτα άλλο.} \\
            }
        }
    }
    & \\ \\

    \theutterance \stepcounter{utterance}  
    & & & \multicolumn{4}{p{0.6\linewidth}}{
        \cellcolor[rgb]{0.9,0.9,0.9}{
            \makecell[{{p{\linewidth}}}]{
                \texttt{\tiny{[GM$\langle$P2]}}
                \texttt{ΕΡΩΤΗΣΗ: Στο πλέγμα σας, η τρίτη γραμμή έχει X μόνο στις τρεις πρώτες στήλες και τα υπόλοιπα δύο κελιά είναι κενά;} \\
            }
        }
    }
     \\ \\

    \theutterance \stepcounter{utterance}  
    & & & \multicolumn{2}{p{0.3\linewidth}}{
        \cellcolor[rgb]{0.9,0.9,0.9}{
            \makecell[{{p{\linewidth}}}]{
                \texttt{\tiny{[GM$|$GM]}}
                \texttt{continue} \\
            }
        }
    }
    & & \\ \\

    \theutterance \stepcounter{utterance}  
    & & \multicolumn{4}{p{0.6\linewidth}}{
        \cellcolor[rgb]{0.9,0.9,0.9}{
            \makecell[{{p{\linewidth}}}]{
                \texttt{\tiny{[P1$\langle$GM]}}
                \texttt{ΑΠΑΝΤΗΣΗ: Η πρώτη γραμμή έχει X μόνο στην πρώτη στήλη, τα υπόλοιπα είναι κενά.} \\
\texttt{ΕΡΩΤΗΣΗ: Στο πλέγμα σας, η τρίτη γραμμή έχει X μόνο στις τρεις πρώτες στήλες και τα υπόλοιπα δύο κελιά είναι κενά;Ξεκινήστε την απάντησή σας με "ΑΠΑΝΤΗΣΗ:" και μην προσθέσετε τίποτα άλλο.} \\
            }
        }
    }
    & \\ \\

    \theutterance \stepcounter{utterance}  
    & \multicolumn{4}{p{0.6\linewidth}}{
        \cellcolor[rgb]{0.9,0.9,0.9}{
            \makecell[{{p{\linewidth}}}]{
                \texttt{\tiny{[P1$\rangle$GM]}}
                \texttt{ΑΠΑΝΤΗΣΗ: Ναι, η τρίτη γραμμή έχει X στις τρεις πρώτες στήλες και τα υπόλοιπα δύο είναι κενά.} \\
            }
        }
    }
    & & \\ \\

    \theutterance \stepcounter{utterance}  
    & & & \multicolumn{2}{p{0.3\linewidth}}{
        \cellcolor[rgb]{0.9,0.9,0.9}{
            \makecell[{{p{\linewidth}}}]{
                \texttt{\tiny{[GM$|$GM]}}
                \texttt{continue} \\
            }
        }
    }
    & & \\ \\

    \theutterance \stepcounter{utterance}  
    & & \multicolumn{4}{p{0.6\linewidth}}{
        \cellcolor[rgb]{0.9,0.9,0.9}{
            \makecell[{{p{\linewidth}}}]{
                \texttt{\tiny{[P1$\langle$GM]}}
                \texttt{Τώρα κάντε μια ερώτηση για να μάθετε νέες πτυχές του πλέγματός μου που μπορεί να διαφέρουν από το δικό σας πλέγμα. Ξεκινήστε με "ΕΡΩΤΗΣΗ:" και μην προσθέσετε τίποτα άλλο.} \\
            }
        }
    }
    & \\ \\

    \theutterance \stepcounter{utterance}  
    & \multicolumn{4}{p{0.6\linewidth}}{
        \cellcolor[rgb]{0.9,0.9,0.9}{
            \makecell[{{p{\linewidth}}}]{
                \texttt{\tiny{[P1$\rangle$GM]}}
                \texttt{ΕΡΩΤΗΣΗ: Η πέμπτη στήλη του πλέγματός σας έχει Χ μόνο στην τελευταία γραμμή και τα υπόλοιπα κελιά της είναι κενά;} \\
            }
        }
    }
    & & \\ \\

    \theutterance \stepcounter{utterance}  
    & & & \multicolumn{2}{p{0.3\linewidth}}{
        \cellcolor[rgb]{0.9,0.9,0.9}{
            \makecell[{{p{\linewidth}}}]{
                \texttt{\tiny{[GM$|$GM]}}
                \texttt{continue} \\
            }
        }
    }
    & & \\ \\

    \theutterance \stepcounter{utterance}  
    & & \multicolumn{4}{p{0.6\linewidth}}{
        \cellcolor[rgb]{0.9,0.9,0.9}{
            \makecell[{{p{\linewidth}}}]{
                \texttt{\tiny{[GM$\rangle$P2]}}
                \texttt{ΑΠΑΝΤΗΣΗ: Ναι, η τρίτη γραμμή έχει X στις τρεις πρώτες στήλες και τα υπόλοιπα δύο είναι κενά.} \\
\texttt{ΕΡΩΤΗΣΗ: Η πέμπτη στήλη του πλέγματός σας έχει Χ μόνο στην τελευταία γραμμή και τα υπόλοιπα κελιά της είναι κενά;Ξεκινήστε την απάντησή σας με "ΑΠΑΝΤΗΣΗ:" και μην προσθέσετε τίποτα άλλο.} \\
            }
        }
    }
    & \\ \\

    \theutterance \stepcounter{utterance}  
    & & & \multicolumn{4}{p{0.6\linewidth}}{
        \cellcolor[rgb]{0.9,0.9,0.9}{
            \makecell[{{p{\linewidth}}}]{
                \texttt{\tiny{[GM$\langle$P2]}}
                \texttt{ΑΠΑΝΤΗΣΗ: Όχι, η πέμπτη στήλη είναι εντελώς κενή σε όλες τις γραμμές.} \\
            }
        }
    }
     \\ \\

    \theutterance \stepcounter{utterance}  
    & & & \multicolumn{2}{p{0.3\linewidth}}{
        \cellcolor[rgb]{0.9,0.9,0.9}{
            \makecell[{{p{\linewidth}}}]{
                \texttt{\tiny{[GM$|$GM]}}
                \texttt{continue} \\
            }
        }
    }
    & & \\ \\

    \theutterance \stepcounter{utterance}  
    & & \multicolumn{4}{p{0.6\linewidth}}{
        \cellcolor[rgb]{0.9,0.9,0.9}{
            \makecell[{{p{\linewidth}}}]{
                \texttt{\tiny{[GM$\rangle$P2]}}
                \texttt{Τώρα πάρε μια απόφαση. Τι πιστεύεις: είναι το δικό σου πλέγμα και το άλλο πλέγμα που περιγράφεται το ίδιο; Γράψε "ΑΠΟΦΑΣΗ: ίδιο πλέγμα" αν πιστεύεις ότι είναι το ίδιο πλέγμα ή "ΑΠΟΦΑΣΗ: διαφορετικά πλέγματα" αν πιστεύεις ότι είναι διαφορετικά πλέγματα. Μην προσθέσεις τίποτα άλλο.} \\
            }
        }
    }
    & \\ \\

    \theutterance \stepcounter{utterance}  
    & & & \multicolumn{4}{p{0.6\linewidth}}{
        \cellcolor[rgb]{0.9,0.9,0.9}{
            \makecell[{{p{\linewidth}}}]{
                \texttt{\tiny{[GM$\langle$P2]}}
                \texttt{ΑΠΟΦΑΣΗ: διαφορετικά πλέγματα} \\
            }
        }
    }
     \\ \\

    \theutterance \stepcounter{utterance}  
    & & & \multicolumn{2}{p{0.3\linewidth}}{
        \cellcolor[rgb]{0.9,0.9,0.9}{
            \makecell[{{p{\linewidth}}}]{
                \texttt{\tiny{[GM$|$GM]}}
                \texttt{continue} \\
            }
        }
    }
    & & \\ \\

    \theutterance \stepcounter{utterance}  
    & & & \multicolumn{2}{p{0.3\linewidth}}{
        \cellcolor[rgb]{0.9,0.9,0.9}{
            \makecell[{{p{\linewidth}}}]{
                \texttt{\tiny{[GM$|$GM]}}
                \texttt{success} \\
            }
        }
    }
    & & \\ \\

    \theutterance \stepcounter{utterance}  
    & & \multicolumn{4}{p{0.6\linewidth}}{
        \cellcolor[rgb]{0.9,0.9,0.9}{
            \makecell[{{p{\linewidth}}}]{
                \texttt{\tiny{[P1$\langle$GM]}}
                \texttt{ΑΠΑΝΤΗΣΗ: Όχι, η πέμπτη στήλη είναι εντελώς κενή σε όλες τις γραμμές.} \\
\texttt{Τώρα πάρε μια απόφαση. Τι πιστεύεις: είναι το δικό σου πλέγμα και το άλλο πλέγμα που περιγράφεται το ίδιο; Γράψε "ΑΠΟΦΑΣΗ: ίδιο πλέγμα" αν πιστεύεις ότι είναι το ίδιο πλέγμα ή "ΑΠΟΦΑΣΗ: διαφορετικά πλέγματα" αν πιστεύεις ότι είναι διαφορετικά πλέγματα. Μην προσθέσεις τίποτα άλλο.} \\
            }
        }
    }
    & \\ \\

    \theutterance \stepcounter{utterance}  
    & \multicolumn{4}{p{0.6\linewidth}}{
        \cellcolor[rgb]{0.9,0.9,0.9}{
            \makecell[{{p{\linewidth}}}]{
                \texttt{\tiny{[P1$\rangle$GM]}}
                \texttt{ΑΠΟΦΑΣΗ: διαφορετικά πλέγματα} \\
            }
        }
    }
    & & \\ \\

    \theutterance \stepcounter{utterance}  
    & & & \multicolumn{2}{p{0.3\linewidth}}{
        \cellcolor[rgb]{0.9,0.9,0.9}{
            \makecell[{{p{\linewidth}}}]{
                \texttt{\tiny{[GM$|$GM]}}
                \texttt{continue} \\
            }
        }
    }
    & & \\ \\

    \theutterance \stepcounter{utterance}  
    & & & \multicolumn{2}{p{0.3\linewidth}}{
        \cellcolor[rgb]{0.9,0.9,0.9}{
            \makecell[{{p{\linewidth}}}]{
                \texttt{\tiny{[GM$|$GM]}}
                \texttt{success} \\
            }
        }
    }
    & & \\ \\

\end{supertabular}
}

\end{document}
